\begin{figure}[htbp]
\centering
\begin{tikzpicture}[scale=1.4]
    % ========== ROC for F(s) = 1/(s+a) ==========
    % Define a = 2 for this example

    % ROC shaded region (Re(s) > -a = -2)
    \fill[UofSCHorseshoe!20] (-2,-2.5) rectangle (3,2.5);

    % Non-convergent region (hatched)
    \fill[pattern=north east lines, pattern color=UofSCRose!40] (-4,-2.5) rectangle (-2,2.5);

    % Grid
    \draw[help lines, gray!30] (-4,-2.5) grid (3,2.5);

    % Axes
    \draw[thick,->] (-4.3,0) -- (3.5,0) node[right] {$\sigma = \Real(s)$};
    \draw[thick,->] (0,-2.8) -- (0,2.8) node[above] {$\jj\omega = \Imag(s)$};

    % ROC boundary (vertical line at Re(s) = -a)
    \draw[UofSCCongaree, very thick, dashed] (-2,-2.5) -- (-2,2.5);
    \node[UofSCCongaree, font=\small, above] at (-2, 2.6) {ROC boundary};
    \node[UofSCCongaree, font=\small] at (-2, -2.8) {$\Real(s) = -a$};

    % Pole at s = -a
    \node[UofSCGarnet, font=\Large] at (-2, 0) {$\times$};
    \node[UofSCGarnet, font=\small, below left] at (-2.2, -0.15) {$s = -a$};

    % ROC label
    \node[UofSCHorseshoe, font=\bfseries] at (0.5, 1.8) {Region of Convergence};
    \node[UofSCHorseshoe, font=\small] at (0.5, 1.4) {$\Real(s) > -a$};

    % Non-ROC label
    \node[UofSCRose, font=\bfseries] at (-3, 1.8) {No Convergence};
    \node[UofSCRose, font=\small] at (-3, 1.4) {$\Real(s) < -a$};

    % Integral convergence annotation
    \draw[UofSCAtlantic, thick, ->] (1.5, -1.5) -- (1.5, -0.8);
    \node[UofSCAtlantic, font=\footnotesize, align=center] at (1.5, -2) {Integral\\converges};

    % Transfer function box
    \node[draw=UofSCBlack, fill=white, sharp corners, inner sep=8pt, font=\normalsize] at (0, -2.2) {
        $F(s) = \dfrac{1}{s + a} = \Laplace\{e^{-at} u(t)\}$
    };

    % Origin marker
    \node[below right, font=\small] at (0.05, -0.05) {$0$};

    % Tick marks
    \foreach \x in {-3,-1,1,2} {
        \draw[thick] (\x, 0.08) -- (\x, -0.08);
        \node[below, font=\tiny] at (\x, -0.15) {$\x$};
    }
    \node[below, font=\tiny] at (-2, -0.15) {$-a$};
\end{tikzpicture}

\vspace{0.8cm}

% Explanation of ROC and causality
\begin{tikzpicture}[scale=0.85]
    % Left: Causal signal (ROC to right of pole)
    \begin{scope}[xshift=-5.5cm]
        \node[font=\bfseries, UofSCAtlantic] at (1.2, 2.5) {Causal Signal};

        % s-plane mini diagram
        \fill[UofSCHorseshoe!20] (-0.5,-1.5) rectangle (2.5,1.5);
        \draw[->] (-1,0) -- (2.8,0) node[right, font=\tiny] {$\sigma$};
        \draw[->] (0,-1.8) -- (0,1.8) node[above, font=\tiny] {$\jj\omega$};
        \node[UofSCGarnet, font=\normalsize] at (-0.5, 0) {$\times$};
        \draw[UofSCCongaree, thick, dashed] (-0.5,-1.5) -- (-0.5,1.5);
        \node[font=\tiny] at (1, 0.8) {ROC};
        \draw[thick,->] (0.3,0) -- (2,0);
        \node[font=\tiny, right] at (2, 0.3) {$\to +\infty$};

        % Time domain plot
        \begin{scope}[yshift=-3.5cm]
            \draw[->] (-0.5,0) -- (2.8,0) node[right, font=\tiny] {$t$};
            \draw[->] (0,-0.3) -- (0,1.5) node[above, font=\tiny] {$f(t)$};
            \draw[UofSCHorseshoe, very thick, domain=0:2.5, samples=50] plot (\x, {1.2*exp(-1.2*\x)});
            \draw[UofSCGarnet, thick] (-0.5,0) -- (0,0);
            \node[font=\tiny] at (1.2, -0.5) {$f(t) = e^{-at}u(t)$};
            \node[font=\tiny] at (1.2, -0.9) {$f(t) = 0$ for $t < 0$};
        \end{scope}
    \end{scope}

    % Right: Anti-causal signal (ROC to left of pole)
    \begin{scope}[xshift=5.5cm]
        \node[font=\bfseries, UofSCRose] at (1.2, 2.5) {Anti-Causal Signal};

        % s-plane mini diagram
        \fill[UofSCHorseshoe!20] (-2.5,-1.5) rectangle (0.5,1.5);
        \draw[->] (-2.8,0) -- (1,0) node[right, font=\tiny] {$\sigma$};
        \draw[->] (0,-1.8) -- (0,1.8) node[above, font=\tiny] {$\jj\omega$};
        \node[UofSCGarnet, font=\normalsize] at (0.5, 0) {$\times$};
        \draw[UofSCCongaree, thick, dashed] (0.5,-1.5) -- (0.5,1.5);
        \node[font=\tiny] at (-1, 0.8) {ROC};
        \draw[thick,->] (-0.3,0) -- (-2,0);
        \node[font=\tiny, left] at (-2, 0.3) {$-\infty \leftarrow$};

        % Time domain plot
        \begin{scope}[yshift=-3.5cm]
            \draw[->] (-2.8,0) -- (0.5,0) node[right, font=\tiny] {$t$};
            \draw[->] (0,-0.3) -- (0,1.5) node[above, font=\tiny] {$f(t)$};
            \draw[UofSCRose, very thick, domain=-2.5:0, samples=50] plot (\x, {1.2*exp(1.2*\x)});
            \draw[UofSCGarnet, thick] (0,0) -- (0.5,0);
            \node[font=\tiny] at (-1.2, -0.5) {$f(t) = -e^{at}u(-t)$};
            \node[font=\tiny] at (-1.2, -0.9) {$f(t) = 0$ for $t > 0$};
        \end{scope}
    \end{scope}

    % Center annotation
    \node[draw=UofSCBlack, fill=UofSCSandstorm, sharp corners, inner sep=6pt, font=\footnotesize, align=center] at (0, 0) {
        Same $F(s)$\\
        Different ROC\\
        $\Downarrow$\\
        Different $f(t)$
    };
\end{tikzpicture}

\caption{Region of Convergence (ROC) for the Laplace transform $F(s) = 1/(s+a)$. The ROC (Horseshoe green shading) is the half-plane $\Real(s) > -a$, to the right of the pole at $s = -a$ (Garnet $\times$). The hatched region shows where the Laplace integral does not converge. The ROC determines whether the inverse transform yields a causal or anti-causal time-domain signal: ROC extending to $+\infty$ corresponds to the causal signal $e^{-at}u(t)$, while ROC extending to $-\infty$ yields the anti-causal signal $-e^{at}u(-t)$. For control systems, we typically assume causal signals, so the ROC extends rightward from the rightmost pole.}
\label{fig:region_of_convergence}
\end{figure}
